\begin{lemma}\label{fill:tate}
At each of $0$ and $\alpha_+$, let $q$ be a local parameter vanishing
at that point. The punctured torus family has a biholomorphic presentation
\begin{equation}
 \begin{gathered}
 X_q\simeq(\CC^*)^2/\langle h,g\rangle,\\
 h(z,v)=(h_1z,h_2v),\qquad g(z,v)=(c_3z,\lambda c_4v),
 \end{gathered}
 \label{fill:tate-period}
\end{equation}
where
\[
 (h_1,h_2)=
 \begin{cases}
 (c_1q^2,c_2q),&\text{at }0,\\
 (c_1q,c_2),&\text{at }\alpha_+,
 \end{cases}
\]
and $c_i=c_i(q)\in\mathbb{C}^*+O(q)$ are holomorphic units.

For the lift of $F$ chosen in the proof, $g=Fh^3$ at $0$ and
$g=F$ at $\alpha_+$. The exponent vector of $g$ is $(0,0)$;
that of $h$ is $(2,1)$ at $0$ and $(1,0)$ at $\alpha_+$.
Both nonzero vectors are primitive.
\end{lemma}
